\scanpage{7}
\margin{18}

(7.13.10)\quad
\diagram{\begin{tikzcd}
\operatorname{Gr}^i(X_{\bar t})_{\mathbf Q}\arrow[r,"\mathrm{cl}"]\arrow[d] & H^{2i}(X_{\bar t},\mathbf Q_\ell(i))\arrow[d,"\sim"]\\
\operatorname{Gr}^i(X_{\bar s})_{\mathbf Q}\arrow[r,"\mathrm{cl}"] & H^{2i}(X_{\bar s},\mathbf Q_\ell(i))
\end{tikzcd}}.

\deleted{7.13.11.} \uline{Remarque\deleted{\unsure{s}}} 7.13.11.\quad La commutativité des deux derniers diagrammes \deleted{commutatifs}, qu'on peut exprimer en disant que la spécialisation des cycles algébriques est compatible avec la formation des classes de cohomologie $\ell$-adiques associées, n'est pas très satisfaisante sous la forme énoncée, puisqu'on \deleted{\unsure{a}} a été obligé de remplacer les coefficients $\mathbf Z_\ell(i)$ par $\mathbf Q_\ell(i)$, c'est à dire de négliger la torsion. Pour avoir un énoncé qui ne néglige pas la torsion, \deleted{on est}\inserted{il semble qu'on sera} obligé de travailler avec l'anneau de Chow, puisque c'est sur ce dernier que ce trouve défini naturellement l'homomorphisme $\mathrm{cl}$ comme homomorphisme à valeurs dans $H^{2*}(-,\mathbf Z_\ell(*))$, \deleted{\unsure{Mais heureusement,}} comparer (XIV 4.1). Malheureusement, on n'a pas développé pour le moment une théorie de la spécialisation pour l'anneau de Chow\deleted{\unsure{,}} (cf. 7.14), qui permettrait d'écrire \deleted{\unsure{un diagramme}}\inserted{des carrés} \unsure{aommutatifs|commutatifs} analogues à (7.13.9) et (7.13.10) en termes des anneaux de Chow et de la cohomologie à coefficients $\mathbf Z_\ell(*)$. Tout ce qu'on peut dire\inserted{lorsque $f:X\longrightarrow S$ est lisse et propre}, c'est qu'on a un carré commutatif

(7.13.12)\quad
\diagram{\begin{tikzcd}
Z^i(X_t)\arrow[r,"\mathrm{cl}"]\arrow[d,"\sigma"'] & H^{2i}(X_t,\mathbf Z_\ell(i))\\
Z^i(X_s)\arrow[r,"\mathrm{cl}"] & H^{2i}(X_s,\mathbf Z_\ell(i))\arrow[u]\arrow[u,phantom,"\deleted{\unsure{\downarrow}}" description]
\end{tikzcd}},

\deleted{et un carré} analogue à (7.13.9), et un carré correspondant analogue à (7.13.10) qu'on laisse au\deleted{\unsure{tre}} lecteur\inserted{le soin} d'écrire, où $Z^i$ désigne le groupe des cycles algébriques (modulo rien du tout) de codimension $i$, \deleted{et} où les flèches horizontales sont les homomorphismes ``classe de cohomologie associée'' définie dans (SGA 5 IV), et où $\sigma$ désigne l'homomorphisme de spécialisation des cycles algébriques, associant à un cycle $Z_t$ le cycle $Z_s$, où $Z$ est
\margin{6}
